\section{Introduction}
\label{sec:introduction}

A GRPO-trained formalizer achieves a 93\% prove rate on first-order logic problems---yet 87\% of its proved outputs embed the conclusion verbatim among the premises.
The solver succeeds not because the model has learned to reason, but because the formalization trivially entails any conclusion already appearing as a premise.
This structural shortcut is invisible to execution-based evaluation: prove rate registers success regardless of whether the proof reflects genuine derivation or degenerate circular embedding.

\paragraph{The execution-only blind spot.}
Prove rate---the dominant evaluation metric for neurosymbolic formalizers~\citep{logic-lm, linc, satlm}---conflates faithful reasoning with degenerate shortcuts, a gap especially consequential under RL, where the optimizer exploits any shortcut the reward fails to penalize~\citep{skalse-reward-hacking, reward-hacking-survey}.
Recent work applies GRPO~\citep{deepseekmath} to train LLMs that generate solver-verified formal programs~\citep{mellgren-prolog-grpo, logic-rl} using binary execution feedback; this drives rapid PR gains but provides no signal on structural quality.
Process reward models~\citep{lightman-verify-step, math-shepherd} require expensive annotations and do not transfer to arbitrary solver-backed settings.
A separate line of work documents that LLMs already game logical formalization at inference time~\citep{kim-game-formalization, bondarenko-spec-gaming}; what is missing is an annotation-free \emph{diagnostic} for structural integrity in neurosymbolic outputs.

\paragraph{Component necessity testing.}
We address this gap with \emph{Solver Trace Ablation} (STA), a diagnostic applying \emph{component necessity testing}---leave-one-out (LOO) ablation on solver inputs---to quantify structural integrity (Figure~\ref{fig:sta-overview}).
For each proved formalization, STA removes each component individually and re-runs the solver; the fraction of components whose removal causes proof failure yields the \emph{Intensional Structural Reward} (ISR), with $\text{ISR}{=}1$ indicating full necessity and $\text{ISR}{=}0$ complete redundancy.
LOO ablation itself is standard~\citep{demillo1978, jia-harman-survey}; novelty lies in its application as a GRPO-hacking diagnostic and in the empirical phenomena it surfaces.
ISR captures structural rather than semantic integrity (downstream NLI: $\rho{=}{-}0.049$, $p{=}0.187$).

\paragraph{ISR as diagnostic.}
As a post-hoc diagnostic on a GRPO formalizer trained with FOLIO~\citep{folio}, ISR reveals a structural collapse trajectory invisible to PR: $\text{ISR}^{*}{:}\,0.642\to 0.132\to 0.003$ over 100--1000 steps (99.5\% degradation) while PR rises $32\%{\to}96\%$.
Three seeds diverge radically (PR $\sigma{=}47$\,pp), each converging to a qualitatively different mode: conclusion embedding, repetition collapse, and tautological collapse.
An SFT baseline on identical data reaches $\text{PR}{=}13.0\%$ with 92.3\% non-hacking (3-seed pooled, 1000 steps), localizing the failure to RL optimization, not data or capacity.

\paragraph{Meta-level Goodhart.}
Repurposed as a training signal, $r = r_{\text{exec}} \cdot (1 + \beta\,\text{ISR})$, ISR is gamed at its own definition---\emph{meta-level Goodhart's Law}~\citep{manheim-goodhart-taxonomy, karwowski-goodhart-rl, gao-reward-overoptimization}.
At $\beta{=}0.1$, 70.6\% of proved samples exhibit conclusion duplication (H7)---versus 13.7\% under matched random (Fisher $p{<}10^{-38}$; random H7 is highly seed-dependent), showing ISR amplifies tautological collapse by ${\sim}5{\times}$; the one seed with a non-trivial residual reduces to $1.06\times$ random after controlling for a premise-count confound.
At $\beta{=}0.3$, both ISR and random produce extreme cross-seed instability (PR $66{\pm}42\%$ and $54{\pm}35\%$, respectively), confirming that high-dose proxy weighting is harmful regardless of signal content; ISR seeds diverge into three distinct pathways (reward collapse, shortcut dominance, pure metric gaming) while random converges uniformly to H1.
The effect transfers and intensifies in Prolog: 3/3 seeds yield 0\% non-hacking via distinct mechanisms, including \emph{query-as-fact gaming} that bypasses text-match classifiers.
ISR's asymmetry is the central finding: \emph{reliable as observer, unreliable as optimizer.}

\paragraph{Contributions.}
Across 6 reward conditions, 3 seeds, 2 logic types, and up to 1000 steps, we contribute: (\textbf{C1}) ISR as a diagnostic for GRPO reward hacking, with a 99.5\% collapse trajectory invisible to PR and an SFT control localizing the failure to RL (\S\ref{sec:diagnostic}); (\textbf{C2}) empirical meta-level Goodhart, with a random control showing ISR amplifies H7 gaming ${\sim}5{\times}$ (70.6\% vs.\ 13.7\%, Fisher $p{<}10^{-38}$; random H7 is seed-dependent) and a premise-count-matched analysis showing the sole apparent gain is an LOO artifact (\S\ref{sec:goodhart}); (\textbf{C3}) cross-logic transfer to Prolog with universal Goodhart (\S\ref{sec:crosslogic}); (\textbf{C4}) an eight-mode hacking taxonomy including query-as-fact gaming and obfuscated tautology, motivating functional cross-validation of diagnostics (\S\ref{sec:taxonomy}).
These are cautionary results for any work treating a structural proxy as a training signal.
